\scanpage{38}\mspage{44}
Un prédérivateur est appelé un \uline{dérivateur à gauche} s'il satisfait Der 1\inserted{ et \unsure{Der|Dér} 2}, et de plus \deleted{\unsure{l'axiome suivant}} les axiomes suivants (dont voici le premier).

% slanted
\margin{Existence des images directes (stable par induction\unsure{.}}
\uline{Der 3} (g) : Pour toute $f:I\to J$ dans Diag, le foncteur $f^*:\mathbb D(J)\to\mathbb D(I)$ admet un \uline{adjoint à droite} $f_*$.

On appelle $f_*$ le foncteur image directe \inserted{(cohomologique)} pour $f$. L'\unsure{objet|obj} $f_*(F)\inserted{{}\in\mathbb D(J)}$ (\inserted{pour }$F\in\mathbb D(I)$) s'appellera aussi \deleted{de} \inserted{l'objet} ``\uline{cohomologie relative}'' de $I$ sur $J$ (par $f$), à \unsure{coëfficients|coéfficients} dans $F$. On a des iso. \unsure{can|can.}
\[
 (gf)_*\simeq g_*f_*
\]

\uline{Cor.} Les catégories $\mathbb D(I)$ ont des produits finis, les foncteurs $f^*$ et $f_*$ y commutent.

Utilise Der 1 \deleted{\unsure{et D}} pour $I=\varnothing$ et $\inserted{I={}}\unsure{.}e\amalg e$, et Der 3 pour $I\xrightarrow{f}e$.

On a un \deleted{\unsure{foncteur}} \inserted{axiome} dual

\uline{Der 3} (d) : existence des adj. à g. $f_!$ des $f^*$

le foncteur $f_!$ s'appelle ``image directe homologique'' pour $f$, et on a une notion d'\uline{homologie relative} de $I$ sur $J$ (par $f$), à coeff. dans $F$.

Par la suite, la mention (g) ou (d) à côté de la désignation Der $n$ d'un axiome, signifie qu'il s'applique aux dériv. à g., resp. \unsure{à} dr. De \unsure{=|même} pour les énoncés, définitions etc.
% archive sheet number: 37
